\chapter{Discussão}\label{cap:discussao}

\section{Integração entre HVSR e MTD}

Na posição do OBS, o topo do MTD Búzios foi estimado em 83,35~m abaixo do fundo marinho. A isócrona local de 51,14~ms, convertida com a velocidade representativa de 1.549,7~m/s, resultou em uma espessura de 39,62~m e em uma profundidade da base de 122,97~m abaixo do fundo marinho. Para 29/02/2020, considerando a frequência central do HVSR e as seis combinações de parâmetros de velocidade, as profundidades variaram de 90,82 a 128,57~m, com ponto médio de 109,69~m e incerteza absoluta pelo envelope de 18,88~m. Assim, a interface inferida pelo HVSR encontra-se dentro do intervalo vertical ocupado pelo depósito e apresenta maior proximidade com sua base do que com o topo.

Essa correspondência sugere que o contraste de propriedades responsável pela resposta fundamental do HVSR possa estar relacionado à porção basal do primeiro evento do MTC Búzios ou à transição entre o depósito e os sedimentos subjacentes. No entanto, essa associação deve ser interpretada com cautela, uma vez que o HVSR não identifica diretamente uma interface estratigráfica específica, mas responde a contrastes nas propriedades elásticas do meio. Ainda assim, a compatibilidade entre a profundidade estimada pelo HVSR e a base sísmica do depósito constitui uma evidência independente que reforça a possibilidade de que a estrutura interna ou basal do MTD exerça influência significativa na resposta geofísica observada.

\section{Limitações e incertezas}

A interpretação integrada está condicionada por limitações metodológicas e espaciais. O HVSR foi obtido em um único ponto do fundo oceânico e, por isso, não representa diretamente toda a variabilidade lateral do MTC Búzios. Além disso, a conversão de frequência em profundidade depende do modelo unidimensional de quarto de comprimento de onda, das velocidades de onda S estimadas indiretamente a partir de $V_p$ e do coeficiente de Poisson, bem como da hipótese de que o pico fundamental corresponda a um contraste de impedância dominante. A sobreposição dos intervalos obtidos nas duas datas sustenta a estabilidade da ordem de grandeza da estimativa, mas não elimina essas fontes de incerteza.

Também existe uma diferença de escala e finalidade entre os modelos de velocidade empregados. A conversão das isócronas sísmicas utiliza velocidades representativas definidas para cada MTD/MTC, enquanto a estimativa por HVSR utiliza velocidades interpoladas no ponto do OBS e convertidas para $V_s$ por cenários de coeficiente de Poisson. A proximidade entre a profundidade estimada pelo HVSR e a base sísmica do depósito deve, portanto, ser tratada como compatibilidade entre métodos independentes, e não como identificação inequívoca da mesma interface. A validação dessa associação requer controle adicional por perfis de velocidade mais próximos, análise da resposta instrumental e oceanográfica, maior número de estações e integração com a resolução vertical da sísmica de reflexão.
